\section{Batched multi-element inclusion proofs}
\label{sec:multiproof}

When proving inclusion of $m > 1$ target indices at once, independent
single-element proofs waste space: two targets whose recursion paths share
a common ancestor block need only reveal that shared checkpoint once, not
once per proof -- exactly the same idea used by Merkle multiproofs.

\subsection{Construction}

\textsc{BuildMultiProof} runs the same recursive traversal as
\textsc{BuildProof} (Section~\ref{sec:algorithm}), but over the set of all
target pairs simultaneously: at each level, a block is only recursed into
if it contains at least one target; a sibling block is only revealed (as
$\mathrm{ending}$ or $\mathrm{prewrap}$, exactly as before) if it contains
\emph{no} target -- if both children contain targets, both are walked and
no reveal is needed at that level at all. The verifier's
\textsc{EvalMultiBlock} mirrors this exactly, computing each shared value
once for every target whose path passes through it.
Algorithms~\ref{alg:build-multi}--\ref{alg:verify-multi} give all three
procedures; they correspond line by line to \texttt{build\_multiproof}
and \texttt{verify\_multiproof} in the reference implementation.
Throughout, $P$ is the set of distinct target \emph{pair starts}, and for
a set $S$ of pair starts inside $\mathrm{block}(a,k)$ we write $S_L = \{s
\in S : s < a + 2^{k-1}\}$ and $S_R = S \setminus S_L$.

\begin{algorithm}[H]
\caption{$\textsc{BuildMultiProof}(T)$ -- prover}
\label{alg:build-multi}
\begin{lstlisting}[style=pseudo]
P := { t - (t mod 2) : t in T }                 # distinct pair starts
function MRecurse(a, k, S):        # S: non-empty, pair starts in block(a,k)
    if k = 1: return               # base pair always sent raw
    h := 2^(k-1)
    if S_L and S_R both non-empty:
        MRecurse(a, k-1, S_L)
        MRecurse(a+h, k-1, S_R)    # no reveal at this level
    else if S_R is empty:
        MRecurse(a, k-1, S_L)
        reveal prewrap(a+h, k-1)   # right sibling, as in BuildProof
    else:
        reveal ending(a, k-1)      # left sibling, as in BuildProof
        MRecurse(a+h, k-1, S_R)

h_top := 2^(L-1);  P_L := {s in P : s < h_top};  P_R := P \ P_L
if P_L and P_R both non-empty:
    MRecurse(0, L-1, P_L);  MRecurse(h_top, L-1, P_R)
else if P_R is empty:
    MRecurse(0, L-1, P_L);  reveal prewrap(h_top, L-1)
else:
    reveal ending(0, L-1);  MRecurse(h_top, L-1, P_R)
return ({ s -> (D[s], D[s+1]) : s in P }, reveals in the order emitted)
\end{lstlisting}
\end{algorithm}

\begin{algorithm}[H]
\caption{$\textsc{EvalMultiBlock}(a, k, S, \mathit{ckpt})$ -- verifier}
\label{alg:eval-multi}
\begin{lstlisting}[style=pseudo]
function EvalMultiBlock(a, k, S, ckpt):
    if k = 1:                                 # here S = {a}
        (d0, d1) := pairs[a]
        return C(C(C(ckpt, d0), d1), pad_b(ckpt))
    h := 2^(k-1)
    if S_L and S_R both non-empty:
        left_end  := EvalMultiBlock(a, k-1, S_L, ckpt)
        right_end := EvalMultiBlock(a+h, k-1, S_R, left_end)
                          # right child seeded with the COMPUTED left_end
    else if S_R is empty:
        left_end    := EvalMultiBlock(a, k-1, S_L, ckpt)
        sib_prewrap := read next reveal       # reject if queue empty
        right_end   := C(sib_prewrap, pad_b(left_end))
    else:
        left_end  := read next reveal         # reject if queue empty
        right_end := EvalMultiBlock(a+h, k-1, S_R, left_end)
    return C(right_end, pad_b(ckpt))
\end{lstlisting}
\end{algorithm}

\begin{algorithm}[H]
\caption{$\textsc{VerifyMultiProof}(N, \Delta, \{\mathit{data}_t\}_{t \in T}, \pi)$ -- top-level verifier}
\label{alg:verify-multi}
\begin{lstlisting}[style=pseudo]
if N < 4 or N is not a power of two: return false
if T is empty or any t in T is outside [0, N): return false
(pairs, reveals) := pi                 # reveals: shared FIFO queue
P := { t - (t mod 2) : t in T }        # derived from the verifier's OWN T,
                                       # never from pi (the a0 rule, Sec. 3)
if keys(pairs) != P: return false      # exactly one pair per pair start
for t in T:
    (d0, d1) := pairs[t - (t mod 2)]
    if t is even: expected := d0 else: expected := d1
    if data_t != expected: return false
h_top := 2^(L-1);  P_L := {s in P : s < h_top};  P_R := P \ P_L
if P_L and P_R both non-empty:
    y := EvalMultiBlock(0, L-1, P_L, IV)
    raw_state := EvalMultiBlock(h_top, L-1, P_R, y)
else if P_R is empty:
    y := EvalMultiBlock(0, L-1, P_L, IV)
    sib_prewrap := read next reveal    # reject if queue empty
    raw_state := C(sib_prewrap, pad_b(y))
else:
    r := read next reveal              # reject if queue empty
    raw_state := EvalMultiBlock(h_top, L-1, P_R, r)
if reveals is not empty: return false  # leftover reveals: malformed
return finalize(raw_state, N) = Delta
\end{lstlisting}
\end{algorithm}

Deriving $P$ from the verifier's own $T$, and requiring the proof's pair
set to equal it exactly, is the multi-target form of the rule
Remark~\ref{rem:a0-bug} establishes for single proofs: the recursion never
looks at a pair's low bit, so a verifier that took pair starts, or the
assignment of indices to pair elements, from the proof could be fooled by
the same relabeling.

\subsection{Soundness}
\label{sec:multiproof-soundness}

\textsc{EvalMultiBlock} (Algorithm~\ref{alg:eval-multi}) is a new
verifier composition, not a parameter change to \textsc{EvalBlock}, and
this paper already flags one soundness pitfall that a superficially
similar-looking verifier design invites (Remark~\ref{rem:prewrap-bug}), so
it needs its own argument rather than an appeal to
Section~\ref{sec:security} by analogy alone.

\begin{lemma}[Multiproof soundness reduces to single-element soundness]
\label{lem:multiproof-soundness}
Fix a genuine digest $\Delta$ committing to a \emph{canonically built} file
(the same scope Theorem~\ref{thm:soundness} requires). Under
Assumptions~\ref{ass:freestart} and~\ref{ass:mdaccess}, no adversary can
produce a target set $T$, values $\{\mathit{data}_t\}_{t \in T}$, and a
multiproof $\pi$ accepted by \textsc{VerifyMultiProof} against $\Delta$,
with $\mathit{data}_t \ne D[t]$ for some $t \in T$, without also producing
a free-start collision in $\C$.
\end{lemma}

\begin{proof}
Fix any $t \in T$ with $\mathit{data}_t \ne D[t]$, and let $a_0 = t - (t
\bmod 2)$. Run \textsc{VerifyMultiProof} on $\pi$ and record every
intermediate value. From this record we extract an ordinary
single-element proof $\pi_t = (d_0, d_1, \mathit{reveals}_t)$ for index
$t$: $(d_0,d_1) = \mathit{pairs}[a_0]$, and $\mathit{reveals}_t$ is built
by walking $t$'s own root-to-leaf path in \textsc{BuildProof} order,
taking at each level (the top-level split included) the value that
\textsc{EvalBlock} would read there:
\begin{itemize}
\item If $t$'s sibling at that level contains no other target pair, $\pi$
  revealed it explicitly, at the same point of the same fold; copy that
  reveal.
\item If the sibling is the \emph{left} child and contains targets,
  \textsc{EvalMultiBlock} computed its value $\mathit{left\_end}$
  recursively and seeded $t$'s child with it; use that computed value as
  the reveal \textsc{EvalBlock} reads in its right branch.
\item If the sibling is the \emph{right} child and contains targets,
  \textsc{EvalMultiBlock} computed
  $\mathit{right\_end} = \textsc{EvalMultiBlock}(a{+}h, k{-}1, S_R,
  \mathit{left\_end})$, whose last step is the call $\C(x,
  \mathrm{pad}_b(\mathit{left\_end}))$; use the pre-closing state $x$ as
  the $\mathit{sib\_prewrap}$ reveal. \textsc{EvalBlock} then performs the
  same call, $\C(x, \mathrm{pad}_b(\mathit{left\_end}))$, with the same
  inputs.
\end{itemize}
By induction on block size along $t$'s path, \textsc{EvalBlock} on
$\pi_t$ performs exactly the compression calls \textsc{EvalMultiBlock}
performs along that path, with identical inputs: the base case uses the
same $(\mathit{ckpt}, d_0, d_1)$, and each level up combines the same
target-side value with the same sibling-side value through the same
Equations~\eqref{eq:foldin}--\eqref{eq:own-close}. Both therefore reach
the same raw chain state and the same $\Delta$, and \textsc{Verify}
selects $\mathit{data}_t$ from $(d_0,d_1)$ by the same parity of $t$ that
\textsc{VerifyMultiProof} already checked. So \textsc{Verify} accepts
$(t, \mathit{data}_t, \pi_t)$ with $\mathit{data}_t \ne D[t]$, and
Theorem~\ref{thm:soundness} supplies a free-start collision. The
extraction costs $O(\log N)$ beyond the multiproof verification itself.
\end{proof}

\begin{remark}[Scope: this is soundness against physical data, not binding]
\label{rem:multiproof-scope}
Lemma~\ref{lem:multiproof-soundness}'s conclusion is tied to the
\emph{true, physical} $D[t]$, exactly as Theorem~\ref{thm:soundness}'s is
-- and, exactly as with the single-element case, this does \emph{not}
extend to a possibly non-canonical file. The extraction above still goes
through unchanged for an arbitrary (possibly non-canonical) $\Delta'$: it
only relies on $\pi$ and the honestly-extracted $\pi_t$ reaching the same
$\Delta'$, not on $\Delta'$ committing to a canonical file. What that
weaker setting yields is the multi-target analog of
Theorem~\ref{thm:binding} (\emph{binding}, not soundness against physical
data): for a fixed target set $T$ and digest $\Delta'$, at most one tuple
$\{\mathit{data}_t\}_{t\in T}$ is accepted by \textsc{VerifyMultiProof},
regardless of whether $\Delta'$ commits to a canonical file -- since two
distinct accepted tuples would disagree at some $t \in T$, extract to two
distinct accepted single-element proofs for that $t$, and
Theorem~\ref{thm:binding} already rules that out. It does \emph{not}
follow that an accepted $\mathit{data}_t$ equals the bytes physically
present at $\mathrm{position}(t)$ in a non-canonical $\Delta'$ --
Remark~\ref{rem:layout-forgery} is a standing counterexample to exactly
that stronger claim, and nothing about batching several targets together
repairs it.
\end{remark}

This reduction is why \textsc{EvalMultiBlock}'s two-case rule (reveal
directly vs.\ recurse) needs to be exactly the single-proof rule applied
independently to each shared block, with no shortcut that lets a
multi-target block reuse a value \emph{without} recomputing or revealing it
per level -- precisely the discipline the insecure alternative of
Remark~\ref{rem:prewrap-bug} violates in the single-target case, and the
same tamper-rejection discipline was applied to
the multiproof implementation as well -- tampering the target data, the
claimed data, one shared or individual reveal per trial, relabeling a
claim to the other index of its pair, and adding an extra unclaimed pair,
across many randomized target sets (Section~\ref{sec:intro}), with every
tamper correctly rejected.

\subsection{Measured savings}

We measured average multiproof size (over 12 random trials per row) against
$m$ independent single-element proofs, for $m$ uniformly random distinct
target indices, at several scales.

\begin{table}[htbp]
\centering
\begin{tabular}{@{}rrrrr@{}}
\toprule
$N$ & $m$ & Independent (B) & Multiproof (B) & Saved \\
\midrule
$2^{10}$ & 2  &   704.0 &   597.3 & 15.2\% \\
$2^{10}$ & 8  &  2816.0 &  1845.3 & 34.5\% \\
$2^{10}$ & 32 & 11264.0 &  5306.7 & 52.9\% \\
$2^{10}$ & 64 & 22528.0 &  8645.3 & 61.6\% \\
\midrule
$2^{14}$ & 2  &   960.0 &   874.7 &  8.9\% \\
$2^{14}$ & 8  &  3840.0 &  2816.0 & 26.7\% \\
$2^{14}$ & 32 & 15360.0 &  9434.7 & 38.6\% \\
$2^{14}$ & 64 & 30720.0 & 16634.7 & 45.9\% \\
\midrule
$2^{18}$ & 2  &  1216.0 &  1120.0 &  7.9\% \\
$2^{18}$ & 8  &  4864.0 &  3896.0 & 19.9\% \\
$2^{18}$ & 32 & 19456.0 & 13509.3 & 30.6\% \\
$2^{18}$ & 64 & 38912.0 & 24965.3 & 35.8\% \\
\midrule
$2^{20}$ & 2  &  1344.0 &  1232.0 &  8.3\% \\
$2^{20}$ & 8  &  5376.0 &  4418.7 & 17.8\% \\
$2^{20}$ & 32 & 21504.0 & 15450.7 & 28.1\% \\
$2^{20}$ & 64 & 43008.0 & 28853.3 & 32.9\% \\
\bottomrule
\end{tabular}
\caption{Multiproof space savings for uniformly random target sets, $D=W=32$
bytes. Savings grow with $m$ (more sharing opportunities) and shrink with
$N$ (a fixed $m$ becomes a sparser sample of a larger tree, so paths share
fewer ancestors), ranging from 7.9\% to 61.6\% across the scales tested.}
\label{tab:multiproof-savings}
\end{table}
